%---------------------------------------------------------------------
\subsection{Intervention Studies}

\subsubsection{Study Summary}

Table~\ref{tab:studies} summarises all primary intervention studies.

\begin{longtable}{>{\raggedright\arraybackslash}p{2.6cm}
                  >{\raggedright\arraybackslash}p{1.9cm}
                  >{\raggedright\arraybackslash}p{3.4cm}
                  >{\raggedright\arraybackslash}p{4.0cm}
                  >{\raggedright\arraybackslash}p{1.8cm}}
\caption{Hangboard intervention studies, 2012--2025.}\label{tab:studies}\\
\toprule
\textbf{Study} & \textbf{Design/$n$} & \textbf{Protocol} & \textbf{Main finding} & \textbf{Cite} \\
\midrule
\endfirsthead
\multicolumn{5}{c}{\textit{(continued)}}\\
\toprule
\textbf{Study} & \textbf{Design/$n$} & \textbf{Protocol} & \textbf{Main finding} & \textbf{Cite} \\
\midrule
\endhead
\bottomrule
\endfoot
\bottomrule
\endlastfoot

López-Rivera \& G-B 2012 & 2-arm RCT; $n=16$ elite & 8~wk, 2$\times$/wk: max weight on 18~mm vs.\ minimal edge, same duration & Max-weight arm superior on strength; minimal-edge arm superior on endurance transfer & \cite{LRGB12} \\[3pt]

Medernach et al.\ 2015 & 2-arm RCT; $n=23$ adv.\ male & 4~wk, 3$\times$/wk fingerboard vs.\ bouldering only & FB: $+2.5$~kg grip, $+5$--7~s DH, $+26$~s IFH vs.\ control & \cite{MKL15} \\[3pt]

López-Rivera \& G-B 2019 & 3-arm RCT; $n=26$ advanced & 8~wk, 2$\times$/wk: MaxHangs / IntHangs / Combo & Endurance: IntHangs $+45\%$ (ES~1.0), MaxHangs $+34\%$ (ES~0.6), Combo $+7\%$; MFS: MaxHangs $+28\%$ only & \cite{LRGB19} \\[3pt]

Levernier \& Laffaye 2019 & Controlled; $n=23$ elite & 4~wk, 2$\times$/wk unilateral hangs (half-crimp $+$ slope) & $+8\%$ MFS; $\uparrow$RFD; control unchanged & \cite{LL19} \\[3pt]

Mundry et al.\ 2021 & 3-arm RCT; $n=30$ int--adv & 8~wk, 2--3$\times$/wk: added weight on 20~mm / progressive edge reduction / control & Added weight $>$ control ($p=0.032$, ES~$0.36$); edge reduction $=$ control & \cite{MSA+21} \\[3pt]

Devise et al.\ 2022 & 4-arm RCT; $n=54$ experienced & 4~wk, 2$\times$/wk: F100 / F80 / F60 / control; 12~mm instrumented hold & F100$+$F80: MFS $+12$--14\%; F80 stamina score: 29.5\%$\to$60.4\% (+31~pp); F60: 32.6\%$\to$53.5\% (+21~pp); F80 improves all three outcomes & \cite{DO22} \\[3pt]

Hermans et al.\ 2022 & 2-arm RCT; $n=35$ int--adv & 10~wk, 2$\times$48~min/wk progressive repeaters (BM1000, 23~mm) vs.\ climbing only & Peak force $+17\%$, avg force $+18\%$ (within-group); RFD $+28\%$ within HBT group but between-group difference not significant ($p=0.056$); DH $+12\%$ within HBT & \cite{HO22} \\[3pt]

Dindorf et al.\ 2024 & 2-arm RCT; $n=18$ completers (14M/4F) int--adv & 7~wk: NMES-assisted fingerboard 2$\times$/wk vs.\ conventional fingerboard 3$\times$/wk & NMES non-inferior to conventional training for finger flexor endurance despite ${\sim}33\%$ lower volume; no significant between-group difference & \cite{DO24} \\[3pt]

Gilmore et al.\ 2024 & Retrospective cohort; $n=527$ Crimpd users & Abrahangs ($\geq3\times$/wk, ${\sim}40\%$ MFS) vs.\ Max Hangs ($\geq0.5\times$/wk, 85--95\%) vs.\ Both vs.\ Climbing Only & Abrahangs $+2.5\%$, Max Hangs $+3.2\%$, Both $+5.8\%$, Climbing Only $+0\%$ strength-to-weight; self-selected, self-reported loads & \cite{GKAO24} \\[3pt]

Langer et al.\ 2024 & RCT; $n=31$ female ($n=26$ completed) & 5~wk: off-wall (fingerboard, BM1000, progressive levels) vs.\ on-wall vs.\ control & Trends toward fingerboard superiority on strength (BF$_{10}=1.16$, anecdotal); on-wall preferred for enjoyment; no significant between-group strength differences & \cite{LAS24} \\[3pt]

Ruiz-López et al.\ 2025 & RCT; $n=11$ experienced (of 20 enrolled) & 4~wk: weighted dead-hang (HIMA, 85\%, 14~mm) vs.\ active-pull dynamometer (PIMA) & Both groups significant gains in peak force, RFD$_{200}$, and max hang time ($p \leq 0.05$); no between-group differences; PIMA reduced hand asymmetry & \cite{RLPCBS25} \\[3pt]

\end{longtable}

\subsubsection{Max Hangs vs.\ Intermittent Hangs (Repeaters)}
\label{sec:maxvsrep}

\begin{evidencebox}{MaxHangs vs.\ IntHangs: Intensity-Specificity (López-Rivera 2019) · \evidencebacked · ESS 7}
López-Rivera \& González-Badillo \cite{LRGB19} randomised 26 advanced climbers (${\sim}7c{+}/8a$, mean age 31.7~yr) to MaxHangs (maximal added weight on 18~mm then minimal edge), IntHangs (4--5 sets $\times$ 10~s on / 5~s off on minimal edge), or Combination (4 weeks each in sequence). IntHangs produced the largest endurance gain ($+45\%$, ES~1.0); MaxHangs the largest MFS gain ($+28\%$, ES~0.6); the Combination group improved endurance by only $+7\%$ and gained no additional MFS. Max hangs maximise MFS; repeaters maximise endurance; sequential combination within one mesocycle underperforms both. \textit{Endpoint: dynamometric surrogate (time-to-failure dead hang, grip endurance) — no on-wall performance or climbing grade outcome measured.}
\end{evidencebox}

\begin{bioplausbox}{Concurrent Training Interference: Mechanism for Poor Combination Result · \bioplausible · ESS 2}
The inferior Combination result is consistent with strength--endurance concurrent-training interference but may equally reflect insufficient per-protocol volume when both are compressed into one mesocycle. This mechanism has not been tested directly in climbing studies.
\end{bioplausbox}

Devise et al.\ \cite{DO22} found that 80\% MFS protocols simultaneously improved MFS ($+12$--14\%), stamina (score rising from 29.5\% to 60.4\%, a gain of $+31$ percentage points), and endurance, making high-submaximal loading the pragmatic default for athletes training both qualities in parallel. The 60\% MFS protocol improved stamina by $+21$ percentage points and produced smaller MFS gains, consistent with an intensity--adaptation specificity relationship.

\subsubsection{Added Weight vs.\ Reduced Edge Depth}

\begin{evidencebox}{Added Weight Outperforms Progressive Edge Reduction (Mundry 2021) · \evidencebacked · ESS 7}
Mundry et al.\ \cite{MSA+21} found that progressive added weight on a 20~mm edge significantly outperformed both edge reduction and control on overall grip strength ($p=0.032$, ES~$0.36$, significant on 3 of 7 pinch tests), while progressive edge reduction yielded no significant gains on any measure. The widely held community preference for minimal-edge training (specificity argument) is not supported by the only RCT directly comparing these methods. \textit{Endpoint: dynamometric surrogate (pinch grip strength) — no on-wall performance or climbing grade outcome measured.}
\end{evidencebox}

\subsubsection{Low-Intensity Daily Hangs (``Abrahangs'')}

\begin{bioplausbox}{Observational: Low-Intensity Daily Hangs (Gilmore 2024) · \bioplausible · ESS 5}
Gilmore et al.\ \cite{GKAO24} (retrospective cohort, $n = 527$ self-selected Crimpd app users, self-reported loads) found: Abrahangs-only $+2.5\%$ strength-to-weight (\textit{not} significantly different from climbing-only, $p = 0.12$); Max Hangs-only $+3.2\%$ (significant vs.\ climbing-only, $p < 0.001$); both combined $+5.8\%$ (additive, significant vs.\ either alone). The Abrahangs-only group did not significantly outperform the climbing-only control, making a non-inferiority claim unsupported. The result is consistent with tendon mechanobiology (low-load cyclic stress promotes collagen synthesis; Baar et al.), but the design carries high risk of selection bias, confounding by total climbing volume, and effort-belief effects. Treat as hypothesis-generating only. A practical risk is that non-expert climbers may misjudge ``submaximal,'' leading to cumulative overuse.
\end{bioplausbox}

\subsubsection{Rate of Force Development}

\begin{evidencebox}{Rate of Force Development Response to Hangboarding · \evidencebacked · ESS 7}
Levernier \& Laffaye \cite{LL19} found a significant $+8\%$ RFD increase after 4 weeks of hangboarding in elite climbers. Hermans et al.\ \cite{HO22} reported a $+28\%$ within-group RFD gain in the hangboard group; however, the between-group difference was not statistically significant ($p = 0.056$, ES 0.21). The meta-analysis by Stien et al.\ \cite{SO23} pooled available trials and found a moderate-to-large effect for RFD (SDM = 0.91, 95\%~CI 0.21--1.61). Grip-depth specificity of RFD adaptation is supported by the Hermans finding that gains were larger on the trained 23~mm rung than on a jug.
\end{evidencebox}

\subsubsection{Training Volume and NMES-Assisted Loading (Dindorf et al.\ 2024)}

\begin{evidencebox}{NMES-Assisted Fingerboard Training Non-Inferior at Lower Volume (Dindorf 2024) · \evidencebacked · ESS 7}
Dindorf et al.\ \cite{DO24} randomised 18 intermediate-to-advanced climbers to twice-weekly NMES-assisted fingerboard training or three-times-weekly conventional fingerboard training for 7 weeks. Despite ${\sim}33\%$ lower training volume, the NMES group was non-inferior to the conventional group on finger flexor endurance. No significant between-group differences were observed. The finding suggests that external neuromuscular augmentation can partially compensate for reduced session frequency, with implications for athletes managing scheduling or recovery constraints. Sample size was small ($n = 18$ completers); results should be considered preliminary.
\end{evidencebox}

\subsubsection{Universal Methodological Limitations}

\begin{itemize}[noitemsep, topsep=2pt]
  \item \textbf{Sample size}: $n=11$--54 in all intervention trials (including Dindorf $n=18$); Stien et al.\ \cite{SO23} found 4 of 5 meta-analyzable trials were underpowered; PEDro scores 5--7/10 (median 6).
  \item \textbf{Cohort}: Nearly all participants are intermediate-to-advanced males; Langer et al.\ \cite{LAS24} is the only female-focused trial and found only anecdotal strength advantages.
  \item \textbf{Duration}: 4--10~weeks; tendon structural remodelling ($>\!3$~months) is not captured by any trial.
  \item \textbf{Outcome validity}: All primary outcomes are dynamometric; the Stien meta found no improvements on actual climbing-specific on-wall tests; no RCT demonstrates statistically significant grade improvement.
  \item \textbf{Confounding}: Participants continued normal climbing throughout; causal attribution is limited.
  \item \textbf{Injury data}: No RCT reports injury incidence; Sjöman et al.\ \cite{SGJ23} provides the only injury-risk data (survey-based).
\end{itemize}

\subsubsection*{Risk of Bias (Simplified Per-Study)}

\begin{longtable}{>{\raggedright\arraybackslash}p{2.8cm}
                  >{\raggedright\arraybackslash}p{1.8cm}
                  >{\raggedright\arraybackslash}p{2.0cm}
                  >{\raggedright\arraybackslash}p{2.0cm}
                  >{\raggedright\arraybackslash}p{1.6cm}
                  >{\raggedright\arraybackslash}p{1.5cm}}
\caption{Simplified risk-of-bias summary for intervention trials.}\label{tab:rob}\\
\toprule
\textbf{Study} & \textbf{Allocation conceal.} & \textbf{Participant blinding} & \textbf{Dropout / ITT} & \textbf{PEDro} & \textbf{Overall} \\
\midrule
\endfirsthead
\multicolumn{6}{c}{\textit{(continued)}}\\
\toprule
\textbf{Study} & \textbf{Allocation conceal.} & \textbf{Participant blinding} & \textbf{Dropout / ITT} & \textbf{PEDro} & \textbf{Overall} \\
\midrule
\endhead
\bottomrule
\endfoot
\bottomrule
\endlastfoot

López-Rivera 2012 & Unclear & Not feasible & Not reported & --- & High \\[2pt]
Medernach 2015 & Unclear & Not feasible & Minimal & 5 & High \\[2pt]
Levernier 2019 & Unclear & Not feasible & Not reported & 5 & High \\[2pt]
López-Rivera 2019 & Unclear & Not feasible & 32\% (38→26) & --- & High \\[2pt]
Mundry 2021 & Unclear & Not feasible & 10\% (30→27) & 6 & High \\[2pt]
Devise 2022 & Unclear & Not feasible & Not reported & 7 & Moderate \\[2pt]
Hermans 2022 & Unclear & Not feasible & Not reported & 6 & High \\[2pt]
Dindorf 2024 & Unclear & Not feasible & 18/32 enrolled (44\%) & --- & High \\[2pt]
Langer 2024 & Stated & Not feasible & 16\% (31→26) & --- & Moderate \\[2pt]
Ruiz-López 2025 & Unclear & Not feasible & 45\% (20→11) & --- & High \\[2pt]
Gilmore 2024 & N/A (obs.) & N/A & Self-selected app users & --- & High (obs.) \\[2pt]
\end{longtable}
\noindent\small PEDro scores from Stien et al.\ \cite{SO23}; `---' = not yet scored or trial postdates meta. `Not feasible' = blinding of participants impossible in exercise training trials; assessor blinding not reported in any trial. Dropout/ITT: López-Rivera 2019 analyzed 26 of 38 enrolled (32\% attrition); Ruiz-López 2025 analyzed 11 of 20 enrolled (45\% attrition) — both without ITT analysis.

\subsubsection{Meta-Analytic Evidence}

\begin{evidencebox}{Hangboard Meta-Analysis: Strength and Endurance Gains (Stien 2023) · \evidencebacked · ESS 7}
Stien et al.\ \cite{SO23} conducted a systematic review and meta-analysis (search to January 2021; 5 trials meta-analyzable, $n=110$) of hangboard and resistance training in climbers. Finger strength improved with pooled SDM~= 0.41 (95\%~CI 0.03--0.80); RFD SDM~= 0.91 (95\%~CI 0.21--1.61); forearm endurance SDM~= 1.23 (95\%~CI 0.69--1.77); overall climbing-specific test performance SDM~= 0.57 (95\%~CI 0.24--0.91). PEDro scores ranged 5--7 (median~6/10); four of five trials were underpowered ($\alpha = 0.05$, $\beta = 0.2$). \textbf{No improvement was found in actual climbing-specific on-wall tests.} \textit{Endpoint note: all pooled outcomes are dynamometric surrogate measures (force, endurance, RFD); no included trial demonstrated statistically significant improvement on a functional on-wall climbing test or grade-based outcome.} Included trials: Levernier \& Laffaye \cite{LL19}, Medernach et al.\ \cite{MKL15}, Hermans et al.\ \cite{HO22}, Mundry et al.\ \cite{MSA+21}, and Stien et al.\ (2021, campus board study). Devise et al.\ \cite{DO22}, Gilmore et al.\ \cite{GKAO24}, and 2024--2025 trials postdate the search.
\end{evidencebox}

\subsubsection{Injury and Experience-Dependent Risk (Sjöman et al.\ 2023)}

\begin{evidencebox}{Experience-Dependent Hangboard Injury Risk (Sjöman 2023) · \evidencebacked · ESS 5}
Sjöman et al.\ \cite{SGJ23} administered a web-based survey to the Finnish climbing community ($n=434$). No significant overall association between regular fingerboard training (RFT) and sport-specific pulley injury (SPIIF) was found. Stratified analysis revealed: in climbers with $\geq6$ years experience, RFT was \emph{negatively} associated with SPIIF ($\chi^2[1, n=200]=4.57$, $p=0.03$); in male climbers with $<6$ years who had reached $\geq7a$, RFT was \emph{positively} associated with SPIIF ($\chi^2[1, n=75]=4.61$, $p=0.03$). This is the only study providing quantitative injury-risk data for hangboarding, suggesting a protective effect in experienced climbers and an elevated risk in less-experienced climbers advancing rapidly.
\end{evidencebox}

\subsubsection{Female-Specific Evidence (Langer et al.\ 2024)}

The only trial enrolling exclusively female climbers \cite{LAS24} found only anecdotal evidence for fingerboard superiority over on-wall training on climbing-specific strength (BF$_{10}=1.16$). On-wall training was strongly preferred for enjoyment and intrinsic motivation. All findings in this review otherwise derive from male-advanced cohorts and should not be assumed to generalise to female or beginner populations.
